%% ma: L12-B15-0001
%% lop: 12
%% chuong: 5
%% bai: 15
%% dang: 12.15.1
%% loai: TN4
%% muc: NB
%% dap_an: "C"
%% nguon: "0. ĐỀ THAM KHẢO TN THPT 2025.pdf (D00), Phần I Câu 4"
%% kien_thuc: "Bài 15 (phương trình chính tắc của đường thẳng)"
%% trang_thai: chinh-thuc
\begin{ex}
Trong không gian với hệ trục tọa độ $Oxyz$, phương trình của đường thẳng đi qua điểm $M(1;-3;5)$ và có một vectơ chỉ phương $\vec u(2;-1;1)$ là
\choice{$\dfrac{x-1}{2}=\dfrac{y-3}{-1}=\dfrac{z-5}{1}$}{$\dfrac{x-1}{2}=\dfrac{y-3}{-1}=\dfrac{z+5}{1}$}{\True $\dfrac{x-1}{2}=\dfrac{y+3}{-1}=\dfrac{z-5}{1}$}{$\dfrac{x+1}{2}=\dfrac{y+3}{-1}=\dfrac{z-5}{1}$}
\loigiai{Đường thẳng qua $M(x_0;y_0;z_0)$ có vectơ chỉ phương $(a;b;c)$ có phương trình $\dfrac{x-x_0}{a}=\dfrac{y-y_0}{b}=\dfrac{z-z_0}{c}$, tức là $\dfrac{x-1}{2}=\dfrac{y+3}{-1}=\dfrac{z-5}{1}$.}
\end{ex}
